\section*{Chapitre \ref{ch:g10:the-mole} --- La mole et la masse molaire}

\begin{solution}{exo:g10:the-mole:1}
$M(\ce{O2}) = 2 \times 16.0 = \qty{32.0}{g/mol}$;
$M(\ce{N2}) = 2 \times 14.0 = \qty{28.0}{g/mol}$;
$M(\ce{CH4}) = 12.0 + 4 \times 1.0 = \qty{16.0}{g/mol}$;
$M(\ce{NH3}) = 14.0 + 3 \times 1.0 = \qty{17.0}{g/mol}$;
$M(\ce{CO2}) = 12.0 + 2 \times 16.0 = \qty{44.0}{g/mol}$.
\end{solution}

\begin{solution}{exo:g10:the-mole:2}
Eau~: $n = 36.0 / 18.0 = \qty{2.00}{mol}$. Méthane~:
$n = 4.0 / 16.0 = \qty{0.25}{mol}$.
\end{solution}

\begin{solution}{exo:g10:the-mole:3}
$m(\ce{NaCl}) = 0.25 \times 58.5 = \qty{14.6}{g}$, environ \qty{15}{g}~;
$m(\ce{CO2}) = 2.0 \times 44.0 = \qty{88}{g}$.
\end{solution}

\begin{solution}{exo:g10:the-mole:4}
$N = 0.50 \times \num{6.02e23} = \num{3.0e23}$ \omterm{def:g7:atoms-and-molecules:molecule}{molécules} de dioxygène.
$n = \num{3.0e24} / \qty{6.02e23}{mol^{-1}} = \qty{5.0}{mol}$ de fer.
\end{solution}

\begin{solution}{exo:g10:the-mole:5}
$V = n \times V_m = 0.50 \times 24.1 = \qty{12}{L}$.
$n = V / V_m = 12 / 24.1 = \qty{0.50}{mol}$.
\end{solution}

\begin{solution}{exo:g10:the-mole:6}
La goutte a une masse de \qty{0.050}{g}, donc
$n = 0.050 / 18.0 = \qty{2.8e-3}{mol}$ et
$N = \num{2.8e-3} \times \num{6.02e23} = \num{1.7e21}$ \omterm{def:g7:atoms-and-molecules:molecule}{molécules}.
\end{solution}

\begin{solution}{exo:g10:the-mole:7}
Un \omterm{def:g7:atoms-and-molecules:atom}{atome} d'aluminium (\qty{27.0}{g/mol}) est environ deux fois plus
léger qu'un \omterm{def:g7:atoms-and-molecules:atom}{atome} de fer (\qty{55.8}{g/mol})~: une même masse
d'aluminium contient donc environ deux fois plus d'\omterm{def:g7:atoms-and-molecules:atom}{atomes}. Vérification~:
aluminium $10/27.0 = \qty{0.370}{mol}$, soit \num{2.23e23} \omterm{def:g7:atoms-and-molecules:atom}{atomes}~; fer
$10/55.8 = \qty{0.179}{mol}$, soit \num{1.08e23} \omterm{def:g7:atoms-and-molecules:atom}{atomes}. L'aluminium
l'emporte, d'un facteur 2.1 environ.
\end{solution}

\begin{solution}{exo:g10:the-mole:8}
Un litre correspond à $1/24.1 = \qty{0.0415}{mol}$ de gaz. Dioxyde de
carbone~: $0.0415 \times 44.0 = \qty{1.83}{g}$~; diazote~:
$0.0415 \times 28.0 = \qty{1.16}{g}$. Le dioxyde de carbone est environ
une fois et demie plus dense que l'\omterm{def:g4:air-a-mixture-of-gases:air}{air}, dont la \omterm{def:g10:the-mole:molar-mass}{masse molaire} est proche
de celle du diazote. Le jus en fermentation dégage du dioxyde de
carbone, qui descend et s'accumule près du sol d'une cave fermée, où il
peut asphyxier quiconque y entre.
\end{solution}

\begin{solution}{exo:g10:the-mole:9}
$M(\ce{C6H12O6}) = 6 \times 12.0 + 12 \times 1.0 + 6 \times 16.0 =
\qty{180.0}{g/mol}$. Flèche $\div N_A$~: $n = \num{1.0e22} / \num{6.02e23} =
\qty{1.66e-2}{mol}$~; flèche $\times M$~: $m = \num{1.66e-2} \times 180.0 =
\qty{3.0}{g}$.
\end{solution}

\begin{solution}{exo:g10:the-mole:10}
$M = \qty{342.0}{g/mol}$, donc $n = 5.0 / 342.0 = \qty{1.46e-2}{mol}$~;
$N = \num{1.46e-2} \times \num{6.02e23} = \num{8.8e21}$ \omterm{def:g7:atoms-and-molecules:molecule}{molécules}.
Chacune a 12 \omterm{def:g7:atoms-and-molecules:atom}{atomes} de carbone~: $12 \times \num{8.8e21} = \num{1.1e23}$
\omterm{def:g7:atoms-and-molecules:atom}{atomes} de carbone.
\end{solution}

\begin{solution}{exo:g10:the-mole:11}
Eau~: $1000 / 18.0 = \qty{55.6}{mol}$. Éthanol~:
$M = 2 \times 12.0 + 6 \times 1.0 + 16.0 = \qty{46.0}{g/mol}$, donc
$1000 / 46.0 = \qty{21.7}{mol}$. L'eau a la plus grande \omterm{def:g10:the-mole:amount}{quantité de
matière}, d'un facteur $55.6 / 21.7 \approx 2.6$, qui n'est autre que
$46.0 / 18.0$.
\end{solution}

\begin{solution}{exo:g10:the-mole:12}
Or~: $\qty{3.75}{g}$, $n = 3.75 / 197.0 = \qty{1.90e-2}{mol}$, soit
\num{1.15e22} \omterm{def:g7:atoms-and-molecules:atom}{atomes}. Cuivre~: \qty{1.25}{g},
$n = 1.25 / 63.5 = \qty{1.97e-2}{mol}$, soit \num{1.19e22} \omterm{def:g7:atoms-and-molecules:atom}{atomes}.
Chose surprenante, la bague contient un peu plus d'\omterm{def:g7:atoms-and-molecules:atom}{atomes} de cuivre que
d'\omterm{def:g7:atoms-and-molecules:atom}{atomes} d'or~: un \omterm{def:g7:atoms-and-molecules:atom}{atome} d'or est environ trois fois plus lourd qu'un
\omterm{def:g7:atoms-and-molecules:atom}{atome} de cuivre.
\end{solution}

\begin{solution}{exo:g10:the-mole:13}
\begin{enumerate}
  \item $V = 8.0 \times 6.0 \times 3.0 = \qty{144}{m^3} =
        \qty{1.44e5}{L}$; $n = \num{1.44e5} / 24.1 = \qty{5.98e3}{mol}$;
        $N = \num{5.98e3} \times \num{6.02e23} = \num{3.6e27}$ molecules.
  \item $M = (78 \times 28.0 + 21 \times 32.0 + 1 \times 40.0) / 100 =
        2896 / 100 = \qty{29.0}{g/mol}$.
  \item $m = \num{5.98e3} \times 29.0 = \qty{1.73e5}{g}$, environ
        \qty{173}{kg}~: l'\omterm{def:g4:air-a-mixture-of-gases:air}{air} d'une salle de classe pèse autant que deux
        ou trois adultes.
\end{enumerate}
\end{solution}

\begin{solution}{exo:g10:the-mole:14}
\begin{enumerate}
  \item $n = 1000 / 28.1 = \qty{35.6}{mol}$, donc
        $N = 35.6 \times \num{6.02e23} = \num{2.14e25}$ \omterm{def:g7:atoms-and-molecules:atom}{atomes}.
  \item $m = \qty{1.000}{kg} / \num{2.14e25} = \qty{4.67e-26}{kg}$.
  \item La masse de la sphère et la \omterm{def:g10:the-mole:molar-mass}{masse molaire} donnent sa \omterm{def:g10:the-mole:amount}{quantité de
        matière}, $n = m / M$. Si l'on compte indépendamment les \omterm{def:g7:atoms-and-molecules:atom}{atomes}
        $N$ (à partir du volume de la sphère et de l'espacement des
        \omterm{def:g7:atoms-and-molecules:atom}{atomes} dans le cristal), le rapport $N / n$ est la \omterm{def:g10:the-mole:avogadro}{constante
        d'Avogadro}.
\end{enumerate}
\end{solution}

\begin{solution}{exo:g10:the-mole:15}
$M = 0.758 \times 35.0 + 0.242 \times 37.0 = 26.53 + 8.95 =
\qty{35.5}{g/mol}$~: la valeur du tableau.
\end{solution}

\begin{solution}{pb:g10:the-mole:1}
\textbf{1.} $M(\ce{H2O}) = 2 \times 1.0 + 16.0 = \qty{18.0}{g/mol}$.

\textbf{2.} $m = 250 \times \qty{1.0}{g} = \qty{250}{g}$.

\textbf{3.} $n = 250 / 18.0 = \qty{13.9}{mol}$.

\textbf{4.} Chaque \omterm{def:g7:atoms-and-molecules:molecule}{molécule} a deux \omterm{def:g7:atoms-and-molecules:atom}{atomes} d'hydrogène et un \omterm{def:g7:atoms-and-molecules:atom}{atome}
d'oxygène~: \qty{27.8}{mol} d'\omterm{def:g7:atoms-and-molecules:atom}{atomes} d'hydrogène, \qty{13.9}{mol}
d'\omterm{def:g7:atoms-and-molecules:atom}{atomes} d'oxygène.

\textbf{5.} $N = 13.9 \times \num{6.02e23} = \num{8.36e24}$ \omterm{def:g7:atoms-and-molecules:molecule}{molécules}.

\textbf{6.} $m = \qty{18.0}{g/mol} / \qty{6.02e23}{mol^{-1}} =
\qty{2.99e-23}{g}$.

\textbf{7.} $\num{8.36e24} \times \qty{2.99e-23}{g} = \qty{250}{g}$, la
masse de l'eau du verre, comme il se doit.

\textbf{8.} $\num{8.36e24} / \num{3.16e7} = \num{2.6e17}$ ans~: un nombre
d'années à dix-sept zéros, bien au-delà de tout comptage possible.

\textbf{9.} $\qty{1.335e9}{km^3} \times \num{e9} = \qty{1.335e18}{m^3}$,
et avec $\qty{1}{m^3} = \qty{1000}{L}$, \qty{1.335e21}{L}.

\textbf{10.} $\num{8.36e24} / \num{1.335e21} = \num{6.26e3}$ \omterm{def:g7:atoms-and-molecules:molecule}{molécules}
marquées par litre.

\textbf{11.} $\num{6.26e3} \times 0.250 = \num{1.57e3}$ \omterm{def:g7:atoms-and-molecules:molecule}{molécules}
marquées dans le second verre.

\textbf{12.} Masse $\num{1.335e21} \times \qty{1.0}{kg} =
\qty{1.335e21}{kg} = \qty{1.335e24}{g}$~; $n = \num{1.335e24} / 18.0 =
\qty{7.42e22}{mol}$.

\textbf{13.} $13.9 / \num{7.42e22} = \num{1.87e-22}$.

\textbf{14.} $\num{1.87e-22} \times \num{8.36e24} = \num{1.57e3}$~: la
même réponse qu'à la question 11, comme il se doit.

\textbf{15.} $\num{1.335e21} / 0.250 = \num{5.34e21}$ verres.

\textbf{16.} $\num{8.36e24} / \num{5.34e21} \approx 1570$. Un verre
contient environ 1600 fois plus de \omterm{def:g7:atoms-and-molecules:molecule}{molécules} que les océans ne
contiennent de verres d'eau~; une fois réparties uniformément, les
\omterm{def:g7:atoms-and-molecules:molecule}{molécules} d'un verre en laissent environ 1600 dans chaque verre d'eau de
mer.

\textbf{17.} $1 / \num{6.26e3} = \qty{1.6e-4}{L} = \qty{0.16}{mL}$,
environ trois gouttes.

\textbf{18.} Le volume des océans est une estimation connue à trois ou
quatre chiffres au mieux, et il a été arrondi~; l'eau de mer n'est pas
de l'eau pure (elle contient des sels)~; le verre ne contient pas
exactement \qty{250}{mL}~; et le \omterm{def:g6:pure-substances-and-mixtures:mixture}{mélange} dans tous les océans ne serait
jamais parfaitement uniforme. Seul l'ordre de grandeur est fiable.

\textbf{19.} Environ 1600 \omterm{def:g7:atoms-and-molecules:molecule}{molécules} du premier verre reviennent dans le
second.
\end{solution}
